\documentclass[12pt,a4paper]{article}

%% ─── Codificação & idioma ───────────────────────────────────────────────────
\usepackage[utf8]{inputenc}
\usepackage[T1]{fontenc}
\usepackage[portuguese]{babel}

%% ─── Tipografia ─────────────────────────────────────────────────────────────
\usepackage{lmodern}
\usepackage{microtype}

%% ─── Matemática ─────────────────────────────────────────────────────────────
\usepackage{amsmath, amssymb}
\usepackage{mathtools}

%% ─── Layout ─────────────────────────────────────────────────────────────────
\usepackage[
  top=2.8cm, bottom=2.8cm,
  left=3.0cm, right=2.5cm,
  headheight=16pt
]{geometry}

%% ─── Cores ──────────────────────────────────────────────────────────────────
\usepackage[dvipsnames,table]{xcolor}
\definecolor{navyblue}{HTML}{1A3A5C}
\definecolor{medblue}{HTML}{2E75B6}
\definecolor{lightblue}{HTML}{D6E4F0}
\definecolor{rowalt}{HTML}{EBF3FB}
\definecolor{formulabg}{HTML}{EEF4FB}
\definecolor{alertbg}{HTML}{FFF8E1}
\definecolor{alertborder}{HTML}{E59C0A}
\definecolor{legalbg}{HTML}{F0FDF4}
\definecolor{legalborder}{HTML}{16A34A}
\definecolor{warnbg}{HTML}{FEF2F2}
\definecolor{warnborder}{HTML}{DC2626}

%% ─── Cabeçalho / rodapé ─────────────────────────────────────────────────────
\usepackage{fancyhdr}
\pagestyle{fancy}
\fancyhf{}
\fancyhead[L]{%
  \footnotesize\color{navyblue}%
  \textbf{ETP} --- Aquisição de Material de Limpeza e Higiene%
}
\fancyhead[R]{%
  \footnotesize\color{navyblue}%
  Borba/AM \textbullet{} Lei n.\textsuperscript{o} 14.133/2021%
}
\fancyfoot[C]{\footnotesize\thepage}
\renewcommand{\headrulewidth}{0.5pt}
\renewcommand{\headrule}{\hbox to\headwidth{\color{medblue}\leaders\hrule height\headrulewidth\hfill}}

%% ─── Caixas coloridas ───────────────────────────────────────────────────────
\usepackage{tcolorbox}
\tcbuselibrary{skins, breakable}

\newtcolorbox{legalbox}[1][Base Legal]{%
  enhanced, breakable,
  colback=legalbg, colframe=legalborder,
  boxrule=0.8pt, arc=4pt,
  left=10pt, right=10pt, top=5pt, bottom=5pt,
  title=\textbf{#1},
  fonttitle=\small\bfseries\color{legalborder},
  attach boxed title to top left={yshift=-2mm, xshift=6mm},
  boxed title style={colback=legalbg, colframe=legalborder, arc=3pt, boxrule=0.6pt}
}

\newtcolorbox{alertbox}[1][Atenção]{%
  enhanced, breakable,
  colback=alertbg, colframe=alertborder,
  boxrule=0.8pt, arc=4pt,
  left=10pt, right=10pt, top=5pt, bottom=5pt,
  title=\textbf{#1},
  fonttitle=\small\bfseries,
  attach boxed title to top left={yshift=-2mm, xshift=6mm},
  boxed title style={colback=alertborder!30, colframe=alertborder, arc=3pt, boxrule=0.6pt}
}

\newtcolorbox{warnbox}[1][Requisito Obrigatório]{%
  enhanced, breakable,
  colback=warnbg, colframe=warnborder,
  boxrule=0.8pt, arc=4pt,
  left=10pt, right=10pt, top=5pt, bottom=5pt,
  title=\textbf{#1},
  fonttitle=\small\bfseries\color{warnborder},
  attach boxed title to top left={yshift=-2mm, xshift=6mm},
  boxed title style={colback=warnbg, colframe=warnborder, arc=3pt, boxrule=0.6pt}
}

\newtcolorbox{formulabox}{%
  enhanced, breakable,
  colback=formulabg, colframe=medblue,
  boxrule=0.7pt, arc=4pt,
  left=14pt, right=14pt, top=2pt, bottom=2pt
}

\newtcolorbox{conclusabox}{%
  enhanced,
  colback=navyblue!8, colframe=navyblue,
  boxrule=1pt, arc=5pt,
  left=12pt, right=12pt, top=8pt, bottom=8pt
}

%% ─── Títulos ────────────────────────────────────────────────────────────────
\usepackage{titlesec}
\titleformat{\section}
  {\color{navyblue}\large\bfseries}
  {\thesection.}{0.6em}{}
  [\color{navyblue}\titlerule]
\titlespacing{\section}{0pt}{20pt}{8pt}

\titleformat{\subsection}
  {\color{medblue}\normalsize\bfseries}
  {\thesubsection}{0.6em}{}
\titlespacing{\subsection}{0pt}{14pt}{5pt}

\titleformat{\subsubsection}
  {\color{navyblue!80}\normalsize\bfseries\itshape}
  {\thesubsubsection}{0.5em}{}
\titlespacing{\subsubsection}{0pt}{10pt}{4pt}

%% ─── Tabelas ────────────────────────────────────────────────────────────────
\usepackage{booktabs}
\usepackage{array}
\usepackage{longtable}
\usepackage{multirow}
\usepackage{tabularx}
\usepackage{colortbl}
\renewcommand{\arraystretch}{1.45}

%% ─── Utilitários ────────────────────────────────────────────────────────────
\usepackage{float}
\usepackage{graphicx}
\usepackage{caption}
\captionsetup{font=small, labelfont=bf, skip=4pt}
\usepackage{enumitem}
\setlist[itemize]{leftmargin=1.8em, itemsep=2pt, topsep=3pt, parsep=0pt}
\setlist[enumerate]{leftmargin=1.8em, itemsep=2pt, topsep=3pt, parsep=0pt}
\usepackage{pdflscape}
\usepackage{setspace}

%% ─── Hiperlinks ─────────────────────────────────────────────────────────────
\usepackage[
  colorlinks=true,
  linkcolor=navyblue,
  urlcolor=medblue,
  pdfborder={0 0 0}
]{hyperref}

%% ─── Numeração de notas de rodapé sem reset por seção ───────────────────────
\usepackage{chngcntr}
\counterwithout{footnote}{section}

%% ─── Macros de conveniência ─────────────────────────────────────────────────
\newcommand{\lei}[1]{\textbf{#1}}
\newcommand{\art}[1]{art.\,#1}
\newcommand{\R}{\text{R\$}\,}
\newcommand{\mq}{\,\text{m}^2}

%% ─── Dados do processo ──────────────────────────────────────────────────────
\newcommand{\etpNum}{ETP-SMS/BORBA-001/2025}
\newcommand{\processoNum}{[NÚMERO DO PROCESSO SEI/E-DOC]}
\newcommand{\modalidade}{Pregão Eletrônico --- Sistema de Registro de Preços}
\newcommand{\objeto}{Aquisição de materiais de limpeza e higiene destinados à Rede de
Atenção à Saúde do Município de Borba/AM, por meio de Ata de Registro de Preços}
\newcommand{\unidadeReq}{Secretaria Municipal de Saúde de Borba/AM}
\newcommand{\servidorResp}{[NOME, CARGO E MATRÍCULA DO SERVIDOR RESPONSÁVEL]}
\newcommand{\dataElab}{Borba/AM, \today}

% ============================================================
\begin{document}
% ============================================================

%% ─────────────────────────────────────────────────────────────
%%  CAPA
%% ─────────────────────────────────────────────────────────────
\begin{titlepage}
  \pagecolor{navyblue}\color{white}
  \setlength{\parindent}{0pt}
  \centering
  \vspace*{1cm}

  {\small\bfseries\color{lightblue} PREFEITURA MUNICIPAL DE BORBA}\\[2pt]
  {\small\color{lightblue!70} SECRETARIA MUNICIPAL DE SAÚDE}\\[0.4cm]

  \textcolor{lightblue!50}{\rule{0.85\linewidth}{0.8pt}}\\[0.6cm]

  {\fontsize{11}{14}\selectfont\bfseries\color{lightblue!80}
    ESTUDO TÉCNICO PRELIMINAR --- ETP}\\[0.3cm]

  {\fontsize{22}{28}\selectfont\bfseries
    Aquisição de Materiais de\\[4pt]
    Limpeza e Higiene}\\[0.25cm]

  {\fontsize{13}{16}\selectfont\color{lightblue!80}
    Ata de Registro de Preços}\\[0.6cm]

  \textcolor{lightblue!50}{\rule{0.85\linewidth}{0.8pt}}\\[1cm]

  \begin{tabular}{@{}r@{\hspace{6pt}}p{0.55\linewidth}@{}}
    \textcolor{lightblue!70}{\small\bfseries N.\textsuperscript{o} do ETP:} &
      \small \etpNum \\[4pt]
    \textcolor{lightblue!70}{\small\bfseries Processo:} &
      \small \processoNum \\[4pt]
    \textcolor{lightblue!70}{\small\bfseries Modalidade:} &
      \small \modalidade \\[4pt]
    \textcolor{lightblue!70}{\small\bfseries Unidade Requisitante:} &
      \small \unidadeReq \\[4pt]
    \textcolor{lightblue!70}{\small\bfseries Responsável pela Elaboração:} &
      \small \servidorResp \\[4pt]
    \textcolor{lightblue!70}{\small\bfseries Elaborado em:} &
      \small \dataElab \\[4pt]
    \textcolor{lightblue!70}{\small\bfseries Base Legal:} &
      \small Lei n.\textsuperscript{o} 14.133/2021 e Decreto n.\textsuperscript{o}
      11.462/2023 (SRP Federal, aplicado por analogia) \\[4pt]
  \end{tabular}

  \vfill
  \textcolor{lightblue!40}{\rule{0.85\linewidth}{0.4pt}}\\[0.3cm]
  {\tiny\color{lightblue!50}
    Documento elaborado em conformidade com o \art{18} e seus incisos da
    Lei n.\textsuperscript{o} 14.133, de 1.\textsuperscript{o} de abril de 2021.}
\end{titlepage}
\nopagecolor
\setlength{\parindent}{1.2cm}

%% ─────────────────────────────────────────────────────────────
%%  IDENTIFICAÇÃO
%% ─────────────────────────────────────────────────────────────
\begin{table}[H]
\centering
\caption*{\textbf{\color{navyblue}FICHA DE IDENTIFICAÇÃO DO ESTUDO TÉCNICO PRELIMINAR}}
\small
\begin{tabular}{>{\bfseries\color{navyblue}}p{4.5cm} p{9.5cm}}
\rowcolor{navyblue!10}\multicolumn{2}{l}{\textbf{\color{navyblue}
  Dados do Processo}} \\
\midrule
N.\textsuperscript{o} do ETP:         & \etpNum \\
\rowcolor{rowalt}
N.\textsuperscript{o} do Processo:    & \processoNum \\
Unidade Requisitante:  & Secretaria Municipal de Saúde de Borba/AM \\
\rowcolor{rowalt}
Unidade Orçamentária:  & [UNIDADE ORÇAMENTÁRIA E FONTE DE RECURSOS] \\
Objeto:                & \objeto \\
\rowcolor{rowalt}
Modalidade Licitatória:& Pregão Eletrônico --- \art{17}, Lei n.\textsuperscript{o} 14.133/2021 \\
Instrumento:           & Ata de Registro de Preços --- \art{82}, Lei n.\textsuperscript{o} 14.133/2021 \\
\rowcolor{rowalt}
Valor Estimado Global: & R\$\,2.000.000,00 (dois milhões de reais) \\
Vigência da ARP:       & 12 (doze) meses, prorrogável por igual período
                         (\art{84}, \S\,1.\textsuperscript{o}) \\
\rowcolor{rowalt}
Elaborado por:         & \servidorResp \\
Data de Elaboração:    & \dataElab \\
\bottomrule
\end{tabular}
\end{table}

\vspace{0.5cm}

%% ─────────────────────────────────────────────────────────────
%%  SUMÁRIO
%% ─────────────────────────────────────────────────────────────
\tableofcontents
\clearpage

% ============================================================
\section{Descrição da Necessidade da Contratação}
% ============================================================

\begin{legalbox}[Base Legal --- \art{18}, I, Lei n.\textsuperscript{o} 14.133/2021]
\small
A contratação deve ser precedida e instruída com ETP, o qual deverá evidenciar o
\textbf{problema a ser resolvido} e sua melhor solução, demonstrando a necessidade
da contratação, fundamentada no planejamento estratégico ou no plano de contratações
anual (PCA), sempre que elaborado, e nas necessidades da sociedade.
\end{legalbox}

\subsection{Contextualização e Identificação do Problema}

O Município de Borba, situado no estado do Amazonas, possui uma população estimada de
\textbf{34.869 habitantes} (IBGE, 2025) e conta com uma rede de atenção à saúde composta
por \textbf{10 unidades assistenciais}: 9 Unidades Básicas de Saúde (UBS) e o
\textbf{Hospital Municipal Vó Mundoca}, único serviço de internação do território
municipal, com 40 leitos operacionais.

A manutenção da \textbf{salubridade ambiental} das unidades de saúde é obrigação legal
irrenunciável, consubstanciada na Resolução RDC ANVISA n.\textsuperscript{o}~216/2004,
RDC n.\textsuperscript{o}~15/2012, na NR-32 (MTE) e no Programa Nacional de Segurança
do Paciente (PNSP/MS, 2013). A ausência ou insuficiência de materiais de limpeza e
higiene acarreta:

\begin{itemize}
  \item Aumento das taxas de \textbf{Infecções Relacionadas à Assistência à Saúde
    (IRAS)}, notoriamente o principal evento adverso evitável em serviços hospitalares;
  \item Descumprimento das exigências sanitárias, com risco de \textbf{interdição
    administrativa} das unidades pela Vigilância Sanitária Municipal e Estadual;
  \item Exposição do trabalhador da saúde a agentes biológicos, em violação à
    \textbf{NR-32} e à responsabilidade civil/trabalhista do Município;
  \item Impacto direto na saúde da população de 34,8 mil pessoas atendidas pela rede.
\end{itemize}

\subsection{Identificação da Área Técnica Demandante}

A área demandante é a \textbf{Secretaria Municipal de Saúde de Borba/AM}, por intermédio
da Gerência de Infraestrutura e Serviços de Saúde, responsável pela gestão dos insumos
não farmacológicos da rede assistencial, conforme organograma institucional vigente.

\subsection{Vinculação ao Planejamento Institucional}

A presente contratação encontra-se prevista no \textbf{Plano de Contratações Anual (PCA)
--- Exercício 2025}, conforme \art{12}, VII, da Lei n.\textsuperscript{o}~14.133/2021,
e está alinhada ao seguinte objetivo estratégico do Plano Municipal de Saúde (PMS
2022--2025): \textit{``Garantir as condições estruturais mínimas para o funcionamento
adequado da rede de atenção à saúde, com foco na segurança do paciente e na prevenção
de doenças infecciosas''.}

% ============================================================
\section{Estimativa das Quantidades a Contratar}
% ============================================================

\begin{legalbox}[Base Legal --- \art{18}, II, Lei n.\textsuperscript{o} 14.133/2021]
\small
O ETP deve conter a estimativa das \textbf{quantidades para a contratação}, acompanhada das memórias de cálculo e dos documentos que lhe dão suporte, considerando a interdependência com outras contratações.
\end{legalbox}

\subsection{Modelagem Matemática da Demanda}

\subsubsection{Área Total de Higienização ($A_H$)}

A área total sob gestão da intendência é a soma das áreas de todas as unidades da rede:

\begin{formulabox}
\[
  A_H \;=\; \sum_{i=1}^{9}\,\text{UBS}_i \;+\; \text{Area}_{\text{Hosp}}
  \;=\; (9 \times 350) + 2.000 \;=\; \mathbf{5.150\,\text{m}^2}
\]
\end{formulabox}

\noindent A subclassificação entre áreas críticas e não críticas, conforme RDC
ANVISA n.\textsuperscript{o}~15/2012, é:

\begin{table}[H]
\centering
\caption{Subclassificação das áreas de higienização da rede de saúde de Borba/AM}
\label{tab:areas}
\small\rowcolors{2}{rowalt}{white}
\begin{tabular}{p{3.8cm}>{\centering\arraybackslash}p{2.2cm}>{\centering\arraybackslash}p{2.2cm}p{5.4cm}}
  \rowcolor{navyblue}
  \color{white}\textbf{Tipo de Área} &
  \color{white}\textbf{Área (m²)} &
  \color{white}\textbf{\% do Total} &
  \color{white}\textbf{Exemplos / Unidades} \\
  \toprule
  Crítica ($A_C$)           & 2.500 & 48{,}5\% &
    Cirurgia, UTI, parto, urgência, expurgo --- Hospital Vó Mundoca. \\
  Não Crítica ($A_{nc}$)    & 2.650 & 51{,}5\% &
    Corredores, recepção, almoxarifado, salas de espera (UBS e Hospital). \\
  \rowcolor{lightblue}
  \textbf{Total ($A_H$)}    & \textbf{5.150} & \textbf{100\%} & \\
  \bottomrule
\end{tabular}
\end{table}

\subsubsection{Consumo Diário de Saneantes ($Q_d$)}

O consumo diário de cada saneante é calculado pela equação de proporcionalidade de área:

\begin{formulabox}
\[
  Q_d \;=\; \left(\frac{A_C}{p_c} \;+\; \frac{A_{nc}}{p_{nc}}\right)\cdot\delta
  \qquad \text{[litros de concentrado por dia]}
\]
\end{formulabox}

\noindent onde $p_c = 200\,\text{m}^2/\text{dia}$ (produtividade em área crítica),
$p_{nc} = 400\,\text{m}^2/\text{dia}$ (área não crítica) e $\delta$ é o coeficiente
de diluição específico de cada produto.

\subsection{Parâmetros de Estoque --- Adaptados à Logística Fluvial}

\subsubsection{Estoque de Segurança ($ES$)}

A condição geográfica de Borba impõe o chamado \textbf{Atrito Logístico Amazônico}:
dependência exclusiva do transporte fluvial, com prazo de reposição médio de 20 dias e
desvio padrão de $\pm\,4$ dias (período cheia) a $\pm\,6$ dias (período seca). O
cálculo estocástico do Estoque de Segurança é:

\begin{formulabox}
\[
  ES \;=\; Z \cdot \sqrt{TR \cdot \sigma_d^2 \;+\; \bar{d}^{\,2} \cdot \sigma_{TR}^2}
\]
\end{formulabox}

\noindent Para insumos de \textbf{criticidade máxima} (Classe A), adota-se
$Z = 3{,}09$ ($99{,}9\%$ de nível de serviço --- distribuição normal padrão).

\subsubsection{Ponto de Pedido ($PP$)}

\begin{formulabox}
\[
  PP \;=\; \left(CMM \times TR_{\text{meses}}\right) \;+\; ES
\]
\end{formulabox}

\noindent Com $TR_{\text{meses}} = 20/30 \approx 0{,}667$ meses (período cheia) e
$TR_{\text{meses}} = 30/30 = 1{,}000$ mês durante a estiagem do Rio Madeira
(julho--outubro).

\subsubsection{Necessidade Real de Aquisição ($N_R$) --- Ciclo de 12 Meses}

\begin{formulabox}
\[
  N_R \;=\; (CMM \times 12) \;+\; ES \;-\; E_{\text{atual}}
\]
\end{formulabox}

\subsection{Quantitativos Estimados por Item}

\begin{alertbox}[Metodologia de Estimativa]
\small
Os quantitativos abaixo foram obtidos a partir de: (i) série histórica de consumo dos
últimos 12 meses; (ii) parâmetros leito/dia do Hospital Vó Mundoca; (iii) parâmetros
m²/dia por área de higienização. Os valores devem ser ratificados pelo Setor de
Almoxarifado antes da publicação do edital, conforme \art{18}, \S\,1.\textsuperscript{o}.
\end{alertbox}

\begin{longtable}{>{\bfseries}p{0.4cm}p{4.8cm}>{\centering\arraybackslash}p{1.5cm}%
  >{\centering\arraybackslash}p{1.2cm}>{\centering\arraybackslash}p{2cm}%
  >{\centering\arraybackslash}p{2.2cm}}

  \caption{Quantitativos estimados para o ciclo de 12 meses}
  \label{tab:quant} \\

  \rowcolor{navyblue}
  \color{white}\textbf{Cl.} &
  \color{white}\textbf{Descrição do Item} &
  \color{white}\textbf{Unid.} &
  \color{white}\textbf{Qtde./mês} &
  \color{white}\textbf{Qtde. anual} &
  \color{white}\textbf{Preço ref. unit. (R\$)} \\
  \toprule
  \endfirsthead
  \multicolumn{6}{l}{\footnotesize\textit{(continuação)}} \\
  \rowcolor{navyblue}
  \color{white}\textbf{Cl.} &
  \color{white}\textbf{Descrição do Item} &
  \color{white}\textbf{Unid.} &
  \color{white}\textbf{Qtde./mês} &
  \color{white}\textbf{Qtde. anual} &
  \color{white}\textbf{Preço ref. unit. (R\$)} \\
  \toprule
  \endhead
  \bottomrule
  \multicolumn{6}{r}{\footnotesize\textit{(continua na próxima página)}} \\
  \endfoot
  \bottomrule
  \endlastfoot
  \rowcolor{rowalt}
  A & Álcool Etílico 70\% --- frasco 500\,mL & Fr & 1.680 & 20.160 & 7,50 \\
  A & Álcool Etílico 70\% --- galão 5\,L & Gl & 168 & 2.016 & 36,00 \\
  \rowcolor{rowalt}
  A & Detergente Enzimático concentrado 1\,L & Fr & 30 & 360 & 42,94 \\
  A & Hipoclorito de Sódio 1\% --- emb.\ 5\,L & Fr & 280 & 3.360 & 5,80 \\
  \rowcolor{rowalt}
  A & Glutaraldeído 2\% --- frasco 1\,L & Fr & 12 & 144 & 65,00 \\
  A & Sabonete Antisséptico clorexidina 2\% 1\,L & Fr & 90 & 1.080 & 18,50 \\
  \rowcolor{rowalt}
  A & Papel Toalha branco 2 dobras (pacote 1000 fls.) & Pct & 36 & 432 & 22,00 \\
  A & Luva de procedimento latex M (cx 100 un.) & Cx & 80 & 960 & 38,00 \\
  \rowcolor{rowalt}
  A & Avental descartável TNT (cx 50 un.) & Cx & 30 & 360 & 55,00 \\
  B & Saco p/\ lixo infectante vermelho 100\,L (rl 100 un.) & Rl & 60 & 720 & 95,00 \\
  \rowcolor{rowalt}
  B & Saco p/\ lixo comum preto 100\,L (rl 100 un.) & Rl & 40 & 480 & 42,00 \\
  B & Desinfetante quaternário de amônio 5\,L & Fr & 45 & 540 & 28,00 \\
  \rowcolor{rowalt}
  B & Álcool Isopropílico 70\% 1\,L & Fr & 25 & 300 & 19,00 \\
  B & Sabão em pó multiuso 1\,kg & Pct & 50 & 600 & 9,50 \\
  \rowcolor{rowalt}
  C & Detergente líquido neutro 500\,mL & Fr & 60 & 720 & 3,80 \\
  C & Água sanitária 2\,L & Fr & 120 & 1.440 & 5,20 \\
  \rowcolor{rowalt}
  C & Vasssoura nylon profissional & Un & 12 & 144 & 18,00 \\
  C & Rodo 60\,cm cabo alumínio & Un & 10 & 120 & 22,00 \\
  \rowcolor{rowalt}
  C & Pano de chão algodão 55$\times$75\,cm & Un & 30 & 360 & 7,50 \\
  C & Balde plástico 20\,L & Un & 5 & 60 & 25,00 \\
  \rowcolor{rowalt}
  C & Esponja dupla face & Un & 40 & 480 & 3,50 \\
\end{longtable}

\begin{alertbox}[Nota sobre Registro de Preços]
\small
A tabela acima é referencial. A Ata de Registro de Preços, por natureza, não obriga
o Município a adquirir o total estimado (\art{83} da Lei n.\textsuperscript{o}~14.133/2021).
As quantidades representam o teto máximo a ser registrado. A contratação efetiva
ocorrerá por ordens de fornecimento emitidas conforme a necessidade da rede,
respeitado o limite da Ata.
\end{alertbox}

% ============================================================
\section{Levantamento de Mercado e Justificativa da Solução}
% ============================================================

\begin{legalbox}[Base Legal --- \art{18}, III, Lei n.\textsuperscript{o} 14.133/2021]
\small
O ETP deve apresentar o \textbf{levantamento de mercado}, com a prospecção e análise das alternativas de solução existentes, e a justificativa da solução escolhida.
\end{legalbox}

\subsection{Alternativas de Solução Analisadas}

\begin{table}[H]
\centering
\caption{Matriz de análise das alternativas de solução}
\label{tab:alternativas}
\small\rowcolors{2}{rowalt}{white}
\begin{tabular}{>{\bfseries}p{3.5cm}p{4.8cm}p{5.0cm}}
  \rowcolor{navyblue}
  \color{white}\textbf{Alternativa} &
  \color{white}\textbf{Vantagens} &
  \color{white}\textbf{Desvantagens / Inviabilidade} \\
  \toprule
  Contrato de Fornecimento Contínuo &
    Previsibilidade de preço e fornecedor. &
    Exige demanda certa; superdimensionamento ou subdimensionamento são comuns;
    não permite flexibilidade nas quantidades. \\
  Compra Direta (dispensa) &
    Agilidade em situações emergenciais. &
    Limitada pelo teto do \art{75} da Lei n.\textsuperscript{o}~14.133/2021;
    não supre demanda anual da rede. \\
  Consórcio Intermunicipal de Compras &
    Economia de escala; redução de custos de frete. &
    Depende de adesão de outros municípios; ausência de consórcio
    operante na região do Rio Madeira para este objeto. \\
  \rowcolor{lightblue}
  \textbf{Ata de Registro de Preços --- Pregão Eletrônico} &
    \textbf{Flexibilidade nas quantidades; preço registrado por 12 meses;
    risco de mercado mitigado; compras parceladas alinhadas ao TR fluvial;
    amplamente utilizada pelo MS para municípios amazônicos.} &
    Requer planejamento prévio robusto para estimativa das quantidades-teto. \\
  \bottomrule
\end{tabular}
\end{table}

\subsection{Justificativa da Solução Escolhida --- Ata de Registro de Preços}

A \textbf{Ata de Registro de Preços (ARP)} é a solução que melhor atende às
particularidades logísticas e operacionais de Borba, pelos seguintes fundamentos:

\begin{enumerate}
  \item \textbf{Adequação ao} \art{82} \textbf{da Lei n.\textsuperscript{o}~14.133/2021}:
    o fornecimento de materiais de limpeza e higiene é caracterizado pela \textit{``entrega
    parcelada''} ao longo do exercício, o que se amolda perfeitamente ao inciso I do
    \art{82}: \textit{``quando, pelas características do bem ou serviço, houver necessidade
    de contratações frequentes''}.
  \item \textbf{Flexibilidade frente ao ciclo hidrológico}: A logística fluvial impõe
    variação nas frequências de entrega. A ARP permite a emissão de ordens de fornecimento
    em volumes diferentes a cada pedido, calibrados ao prazo de reposição $TR$ vigente
    (cheia vs.\ seca do Rio Madeira).
  \item \textbf{Economia de escala}: O registro do preço por 12 meses protege o Município
    contra a sazonalidade de preços dos insumos, frequentemente majorados em Manaus no
    período de seca, quando a demanda fluvial por abastecimento aumenta.
  \item \textbf{Conformidade com práticas do Ministério da Saúde}: O Programa de
    Qualificação das Ações de Vigilância em Saúde (PAVS) recomenda o uso de ARP para
    insumos de saúde em municípios de difícil acesso.
\end{enumerate}

\subsection{Pesquisa de Mercado}

A pesquisa de preços foi realizada em conformidade com a Instrução Normativa SEGES/ME
n.\textsuperscript{o}~65/2021 (aplicada por analogia ao âmbito municipal), mediante:

\begin{itemize}
  \item Consulta ao \textbf{Painel de Preços do Governo Federal}
    (\url{paineldeprecos.planejamento.gov.br}), com base nas compras homologadas
    nos últimos 12 meses para itens idênticos ou similares;
  \item Consulta ao \textbf{BPS --- Banco de Preços em Saúde} do Ministério da Saúde;
  \item Cotações diretas junto a \textbf{3 (três) fornecedores} do mercado de Manaus,
    obtidas por e-mail, com documentação juntada ao processo;
  \item Análise dos \textbf{contratos anteriores} da Secretaria Municipal de Saúde
    (exercícios 2022--2024), disponíveis no sistema de gestão da unidade.
\end{itemize}

O valor estimado global de \textbf{R\$~2.000.000,00} (dois milhões de reais) resultou
da \textbf{média aritmética} dos preços obtidos nas fontes acima, acrescida dos
\textbf{custos logísticos fluviais} Manaus--Borba, conforme matriz de custo por m²
$C_{m^2}$:

\begin{formulabox}
\[
  C_{m^2} \;=\; \frac{\displaystyle
    \sum_{j}\bigl(P_{\text{Manaus},j} \;+\; F_{\text{Borba},j} \;+\; \text{Perdas}_j\bigr)
    \cdot Q_j}{A_H}
\]
\end{formulabox}

\noindent onde $F_{\text{Borba}}$ representa o frete fluvial médio (8--15\% do
valor CIF Manaus) e Perdas representa perdas de transporte e armazenagem
(1--3\% para líquidos).

% ============================================================
\section{Estimativa do Valor da Contratação}
% ============================================================

\begin{legalbox}[Base Legal --- \art{18}, IV e \art{23}, Lei n.\textsuperscript{o} 14.133/2021]
\small
O ETP deve trazer a estimativa do valor da contratação, por preço de mercado, devidamente justificada nos autos, em conformidade com o \art{23}: o valor estimado deve ser obtido por pesquisa de mercado, nos moldes da normativa do órgão federal competente.
\end{legalbox}

\subsection{Metodologia de Estimativa}

A estimativa do valor contratual adotou os seguintes critérios, em ordem de preferência:

\begin{enumerate}
  \item Preço médio apurado no \textbf{Painel de Preços} (SEGES/MGI) para contratações
    iguais ou similares, com data de homologação nos últimos 12 meses;
  \item Preço do \textbf{BPS/MS} para itens hospitalares;
  \item Média das \textbf{3 cotações diretas} junto a fornecedores de Manaus;
  \item Em casos de divergência superior a 50\% entre as fontes, foi descartado o valor
    discrepante e calculada a média das demais, conforme orientação da IN SEGES
    n.\textsuperscript{o}~65/2021.
\end{enumerate}

\subsection{Valor Estimado Global}

\begin{table}[H]
\centering
\caption{Composição do valor estimado por classe ABC}
\label{tab:valor}
\small\rowcolors{2}{rowalt}{white}
\begin{tabular}{>{\centering\arraybackslash}p{1.8cm}p{5.2cm}%
  >{\centering\arraybackslash}p{2.8cm}>{\centering\arraybackslash}p{2.8cm}}
  \rowcolor{navyblue}
  \color{white}\textbf{Classe} &
  \color{white}\textbf{Grupo de Itens} &
  \color{white}\textbf{Valor Estimado (R\$)} &
  \color{white}\textbf{\% do Total} \\
  \toprule
  A & Saneantes críticos, EPI de limpeza, papel toalha hospitalar & 1.600.000,00 & 80\% \\
  B & Sacos para lixo, desinfetantes, sabonetes antissépticos      &   300.000,00 & 15\% \\
  C & Utensílios de limpeza, água sanitária, detergente neutro     &   100.000,00 &  5\% \\
  \rowcolor{lightblue}
  \multicolumn{2}{c}{\textbf{TOTAL GERAL}}                        & \textbf{2.000.000,00} & \textbf{100\%} \\
  \bottomrule
\end{tabular}
\end{table}

\begin{alertbox}[Frete Fluvial Incluso]
\small
O valor estimado já incorpora o \textbf{custo do frete fluvial Manaus--Borba}
(estimado em 8--15\% do preço FOB Manaus) e as perdas de transporte (1--3\%).
O edital deverá exigir que as propostas sejam apresentadas em preço \textbf{CIF Borba},
vedada a oferta de preço FOB Manaus, sob pena de desclassificação.
\end{alertbox}

% ============================================================
\section{Justificativa para o Parcelamento ou Não da Solução}
% ============================================================

\begin{legalbox}[Base Legal --- \art{18}, V, Lei n.\textsuperscript{o} 14.133/2021]
\small
O ETP deve conter a justificativa para o parcelamento ou não da solução de contratação.
\end{legalbox}

\subsection{Parcelamento em Lotes}

A contratação \textbf{será dividida em lotes}, agrupados por \textbf{afinidade técnica
e compatibilidade de fornecimento}, nos termos do \art{40}, \S\,1.\textsuperscript{o}
da Lei n.\textsuperscript{o}~14.133/2021, que determina o parcelamento sempre que for
tecnicamente viável e economicamente vantajoso:

\begin{table}[H]
\centering
\caption{Estrutura de lotes proposta para o Pregão Eletrônico}
\label{tab:lotes}
\small\rowcolors{2}{rowalt}{white}
\begin{tabular}{>{\centering\arraybackslash}p{1.5cm}p{6cm}%
  >{\centering\arraybackslash}p{2.2cm}>{\centering\arraybackslash}p{3.2cm}}
  \rowcolor{navyblue}
  \color{white}\textbf{Lote} &
  \color{white}\textbf{Descrição} &
  \color{white}\textbf{Classe ABC} &
  \color{white}\textbf{Valor Est.\ (R\$)} \\
  \toprule
  01 & Álcoois e antissépticos para mãos e superfícies   & A & 820.000,00 \\
  02 & Saneantes de uso hospitalar (enzimático, glutaraldeído, hipoclorito) & A & 480.000,00 \\
  03 & EPI de limpeza e materiais descartáveis           & A & 300.000,00 \\
  04 & Sacos para resíduos de saúde e comuns             & B & 180.000,00 \\
  05 & Desinfetantes, sabonetes e complementares         & B & 120.000,00 \\
  06 & Utensílios de limpeza e consumíveis domésticos    & C & 100.000,00 \\
  \rowcolor{lightblue}
  \multicolumn{2}{c}{\textbf{TOTAL}}                    &   & \textbf{2.000.000,00} \\
  \bottomrule
\end{tabular}
\end{table}

\subsection{Justificativa do Agrupamento}

O agrupamento em 6 lotes fundamenta-se em:

\begin{enumerate}
  \item \textbf{Compatibilidade de fornecimento}: os itens de cada lote são
    comercializados pelos mesmos distribuidores, permitindo que empresas especializadas
    participem sem necessidade de constituir consórcios;
  \item \textbf{Viabilidade técnica de fracionamento}: cada lote possui dimensão
    suficiente para atrair competição de mercado, conforme \art{40},
    \S\,1.\textsuperscript{o}, I;
  \item \textbf{Gestão de risco}: a independência dos lotes evita que a inadimplência
    de um fornecedor comprometa toda a cadeia de suprimentos;
  \item \textbf{Não fragmentação artificial}: os lotes não subdividem o objeto de forma
    a burlar os limites de dispensa de licitação (\art{75}), vedação expressa no
    \art{34}, \S\,3.\textsuperscript{o}.
\end{enumerate}

% ============================================================
\section{Contratações Correlatas e Interdependentes}
% ============================================================

\begin{legalbox}[Base Legal --- \art{18}, VI, Lei n.\textsuperscript{o} 14.133/2021]
\small
O ETP deve identificar as contratações correlatas ou interdependentes, de modo a
possibilitar o planejamento conjunto das contratações.
\end{legalbox}

\begin{table}[H]
\centering
\caption{Contratações correlatas e interdependentes}
\label{tab:correlatas}
\small\rowcolors{2}{rowalt}{white}
\begin{tabular}{p{5.2cm}p{4.4cm}p{4.0cm}}
  \rowcolor{navyblue}
  \color{white}\textbf{Contratação} &
  \color{white}\textbf{Natureza da Relação} &
  \color{white}\textbf{Status} \\
  \toprule
  Serviços de limpeza e higienização hospitalar (Hosp.\ Vó Mundoca) &
    Interdependente: os materiais deste ETP são utilizados pelo prestador contratado &
    Contrato vigente --- verificar prazo de renovação \\
  Coleta e transporte de resíduos de saúde (RSS) &
    Correlata: os sacos para lixo infectante (Lote 04) integram o fluxo de RSS &
    Vigente ou a contratar \\
  Aquisição de EPI para trabalhadores da saúde &
    Correlata: parte dos EPI de limpeza podem ser cobertos por contrato de EPI geral &
    Verificar sobreposição de objeto \\
  Manutenção de equipamentos de limpeza & Correlata: lavadoras, aspiradores, enceradeiras &
    Verificar manutenção preventiva vigente \\
  \bottomrule
\end{tabular}
\end{table}

\begin{warnbox}[Vedação à Sobreposição de Objeto]
\small
Antes da publicação do edital, a área técnica deve confirmar que os itens
deste ETP não estão cobertos por outros contratos vigentes, sob pena de
\textbf{duplicidade de despesa} e responsabilização do gestor (\art{148},
Lei n.\textsuperscript{o}~14.133/2021).
\end{warnbox}

% ============================================================
\section{Resultados Pretendidos}
% ============================================================

\begin{legalbox}[Base Legal --- \art{18}, VII, Lei n.\textsuperscript{o} 14.133/2021]
\small
O ETP deve demonstrar os resultados pretendidos em termos de qualidade e efetividade da contratação, com indicadores que possibilitem mensurar o grau de cumprimento dos objetivos.
\end{legalbox}

\subsection{Resultados Esperados}

\begin{enumerate}
  \item \textbf{Cobertura ininterrupta}: garantia de abastecimento de saneantes
    críticos (Classe A) para os 5.150\,m² da rede durante os 12 meses da Ata,
    incluindo o período de estiagem do Rio Madeira;
  \item \textbf{Redução de IRAS}: manutenção das taxas de infecção hospitalar do
    Hospital Vó Mundoca abaixo das metas do PNSP/MS;
  \item \textbf{Conformidade sanitária}: zero não conformidades em auditorias da
    Vigilância Sanitária Municipal e Estadual relacionadas à deficiência de insumos
    de limpeza;
  \item \textbf{Eficiência orçamentária}: custo por m² higienizado $\leq$
    R\$~8,50/m², representando redução de no mínimo 10\% em relação à contratação
    anterior.
\end{enumerate}

\subsection{Indicadores de Desempenho (KPIs)}

\begin{table}[H]
\centering
\caption{Indicadores de desempenho da contratação}
\label{tab:kpis}
\small\rowcolors{2}{rowalt}{white}
\begin{tabular}{p{4.0cm}p{3.8cm}>{\centering\arraybackslash}p{2.4cm}%
  >{\centering\arraybackslash}p{2.4cm}}
  \rowcolor{navyblue}
  \color{white}\textbf{Indicador (KPI)} &
  \color{white}\textbf{Fórmula} &
  \color{white}\textbf{Meta} &
  \color{white}\textbf{Periodicidade} \\
  \toprule
  Taxa de Ruptura de Estoque &
    Dias em falta / Dias do período & $< 0{,}1\%$ & Mensal \\
  Custo por m² Higienizado &
    Custo total / $A_H$ & $\leq$ R\$ 8,50/m² & Trimestral \\
  Aderência ao Ponto de Pedido &
    Pedidos no prazo / Total de pedidos & $> 95\%$ & Mensal \\
  Perda por validade / deterioração &
    Valor descartado / Valor adquirido & $< 2\%$ & Semestral \\
  Cobertura do Estoque de Segurança &
    $E_{\text{atual}} / ES_{\text{calc.}}$ & $> 100\%$ & Quinzenal \\
  \bottomrule
\end{tabular}
\end{table}

% ============================================================
\section{Providências para Adequação do Ambiente Organizacional}
% ============================================================

\begin{legalbox}[Base Legal --- \art{18}, VIII, Lei n.\textsuperscript{o} 14.133/2021]
\small
O ETP deve apontar as providências necessárias à adequação do ambiente organizacional, de infraestrutura e de pessoal para receber o objeto da contratação.
\end{legalbox}

\begin{enumerate}
  \item \textbf{Infraestrutura de armazenagem}: adequação do almoxarifado do Hospital
    Vó Mundoca para recebimento e guarda de produtos químicos (prateleiras a $\geq 50$\,cm
    do solo; ventilação adequada; separação de inflamáveis e corrosivos conforme ABNT
    NBR 14.725 e NBR 7.500);
  \item \textbf{Capacitação}: treinamento da equipe de almoxarifado nos procedimentos
    de recebimento, conferência de registro ANVISA, controle de validade (FEFO ---
    \textit{First Expiry, First Out}) e emissão de ordens de fornecimento;
  \item \textbf{Sistema de controle de estoque}: implantação ou atualização do módulo
    de controle de insumos no sistema de gestão da SMS, alimentando os parâmetros
    $CMM$, $ES$, $PP$ e $N_R$ definidos neste ETP;
  \item \textbf{Designação do Gestor e Fiscal do Contrato}: nomeação formal, com
    publicação em Diário Oficial, de servidor para a função de gestor e outro para
    fiscal da Ata de Registro de Preços, nos termos dos arts.~117 a~120 da Lei
    n.\textsuperscript{o}~14.133/2021;
  \item \textbf{Calendário de pedidos}: elaboração de calendário anual de emissão de
    ordens de fornecimento, alinhado ao ciclo hidrológico do Rio Madeira e aos
    parâmetros de $PP$ calculados neste ETP.
\end{enumerate}

% ============================================================
\section{Contratação por Terceiros --- Impossibilidade de Suprimento Interno}
% ============================================================

\begin{legalbox}[Base Legal --- \art{18}, IX, Lei n.\textsuperscript{o} 14.133/2021]
\small
O ETP deve demonstrar que a necessidade não pode ser suprida pela execução direta por órgãos e entidades da Administração Pública.
\end{legalbox}

O Município de Borba \textbf{não dispõe de capacidade produtiva própria} para a fabricação
de saneantes, insumos químicos ou materiais de higiene. Trata-se de objeto de mercado,
com ampla concorrência no segmento de distribuidores atacadistas de Manaus, cujas
empresas possuem autorização de funcionamento junto à ANVISA e são aptas a participar
de pregão eletrônico. A contratação com terceiros é, portanto, a \textbf{única solução
tecnicamente viável}.

% ============================================================
\section{Análise de Riscos da Contratação}
% ============================================================

\begin{legalbox}[Base Legal --- \art{18}, X e Anexo I da Lei n.\textsuperscript{o} 14.133/2021]
\small
O planejamento da contratação deve contemplar a análise dos riscos que possam comprometer o sucesso da contratação e o alcance dos resultados pretendidos.
\end{legalbox}

\begin{longtable}{>{\centering\arraybackslash}p{0.4cm}p{3.8cm}>{\centering\arraybackslash}p{1.2cm}%
  >{\centering\arraybackslash}p{1.2cm}>{\centering\arraybackslash}p{1.2cm}p{4.5cm}}

  \caption{Matriz de riscos da contratação}
  \label{tab:riscos} \\

  \rowcolor{navyblue}
  \color{white}\textbf{\#} &
  \color{white}\textbf{Risco Identificado} &
  \color{white}\textbf{Probab.} &
  \color{white}\textbf{Impacto} &
  \color{white}\textbf{Nível} &
  \color{white}\textbf{Ação de Mitigação} \\
  \toprule
  \endfirsthead
  \multicolumn{6}{l}{\footnotesize\textit{(continuação)}} \\
  \rowcolor{navyblue}
  \color{white}\textbf{\#} &
  \color{white}\textbf{Risco Identificado} &
  \color{white}\textbf{Probab.} &
  \color{white}\textbf{Impacto} &
  \color{white}\textbf{Nível} &
  \color{white}\textbf{Ação de Mitigação} \\
  \toprule
  \endhead
  \bottomrule
  \endlastfoot
  \rowcolor{rowalt}
  R1 & Atraso ou suspensão de balsa no período de seca do Rio Madeira &
    Alta & Alto & \textbf{\color{red}Crítico} &
    Constituição de ERE (60 dias) para Classe A; PP ajustado a $TR = 30$ dias de
    julho a outubro. \\
  R2 & Fornecedor registrado na ARP sem capacidade de entrega CIF Borba &
    Média & Alto & \textbf{\color{orange}Alto} &
    Exigir comprovação de logística fluvial na fase de habilitação; admitir
    no máximo 2 fornecedores por item. \\
  \rowcolor{rowalt}
  R3 & Perda de validade por superdimensionamento ou armazenagem inadequada &
    Média & Médio & \textbf{\color{orange}Médio} &
    Adotar sistema FEFO; limitar pedidos a 2 meses de consumo; treinar
    almoxarifado. \\
  R4 & Subdimensionamento dos quantitativos-teto na ARP &
    Baixa & Alto & \textbf{\color{orange}Médio} &
    Utilizar série histórica e parâmetros matemáticos deste ETP; rever CMM
    semestralmente. \\
  \rowcolor{rowalt}
  R5 & Aumento de preço superior ao IPCA durante vigência da ARP &
    Média & Baixo & \textbf{\color{blue}Baixo} &
    Cláusula de reequilíbrio econômico-financeiro (\art{124}) com indexador
    INPC ou IPCA. \\
  R6 & Descontinuidade de registro ANVISA durante a vigência da ARP &
    Baixa & Médio & \textbf{\color{blue}Baixo} &
    Edital deve exigir confirmação trimestral do registro; sanção por entrega
    de produto sem registro. \\
  \rowcolor{rowalt}
  R7 & Impugnações ao edital por restrições de habilitação &
    Baixa & Baixo & \textbf{\color{gray}Mínimo} &
    Adequar exigências ao mínimo necessário; submeter minuta do edital à
    PGM antes da publicação. \\
\end{longtable}

% ============================================================
\section{Sustentabilidade e Desenvolvimento Nacional}
% ============================================================

\begin{legalbox}[Base Legal --- \art{18}, VII e \art{11}, IV, Lei n.\textsuperscript{o} 14.133/2021]
\small
A Administração deve considerar, nas contratações, a \textbf{promoção do desenvolvimento
nacional sustentável}, conforme o \art{11}, IV, incluindo critérios de sustentabilidade
ambiental (\art{5} do Decreto n.\textsuperscript{o}~7.746/2012).
\end{legalbox}

As especificações técnicas do Termo de Referência deverão incorporar os seguintes
requisitos de sustentabilidade:

\begin{itemize}
  \item \textbf{Saneantes biodegradáveis}: preferência por produtos com laudo de
    biodegradabilidade emitido por laboratório acreditado pelo INMETRO, nos termos
    da RDC ANVISA n.\textsuperscript{o}~59/2021;
  \item \textbf{Embalagens com menor geração de resíduos}: vedação a embalagens
    unitárias desnecessárias; preferência por refis e embalagens com símbolo
    de reciclabilidade;
  \item \textbf{Microempresas e EPP}: aplicação dos arts.~44 e~45 da LC
    n.\textsuperscript{o}~123/2006 (empate ficto), promovendo o desenvolvimento
    econômico local e regional;
  \item \textbf{Destinação de embalagens}: o Termo de Referência deverá conter
    obrigação do fornecedor de recolher e dar destinação ambientalmente adequada
    às embalagens de saneantes, conforme a Política Nacional de Resíduos Sólidos
    (Lei n.\textsuperscript{o}~12.305/2010).
\end{itemize}

% ============================================================
\section{Modalidade, Tipo de Licitação e Regime de Execução}
% ============================================================

\begin{table}[H]
\centering
\caption{Definição da modalidade licitatória}
\label{tab:modalidade}
\small\rowcolors{2}{rowalt}{white}
\begin{tabular}{>{\bfseries\color{navyblue}}p{4.5cm}p{9.5cm}}
\rowcolor{navyblue!10}
\multicolumn{2}{l}{\textbf{\color{navyblue}Configuração da Licitação}} \\
\midrule
Modalidade:           & Pregão Eletrônico (\art{17} c/c \art{28}, I) \\
\rowcolor{rowalt}
Critério de julgamento: & Menor Preço (\art{33}, I) \\
Modo de disputa:      & Aberto (\art{56}, I) --- lances sucessivos com
                        intervalo mínimo de 30 segundos \\
\rowcolor{rowalt}
Instrumento de Contratação: & Ata de Registro de Preços (\art{82}) \\
Vigência da ARP:      & 12 meses, contados da publicação,
                        prorrogável por igual período (\art{84}, \S\,1.\textsuperscript{o}) \\
\rowcolor{rowalt}
Regime de execução:   & Fornecimento parcelado, mediante ordens de
                        fornecimento emitidas pelo órgão gerenciador \\
Participação:         & Ampla, admitida a participação de ME/EPP com
                        benefícios da LC n.\textsuperscript{o}~123/2006 \\
\rowcolor{rowalt}
Sistema eletrônico:   & ComprasNet / PNCP (Portal Nacional de Contratações
                        Públicas), conforme \art{174} \\
\bottomrule
\end{tabular}
\end{table}

% ============================================================
\section{Requisitos de Habilitação}
% ============================================================

Sem prejuízo do Termo de Referência, o ETP antecipa os seguintes requisitos de
habilitação mínimos, em conformidade com os arts.~62 a~70 da Lei
n.\textsuperscript{o}~14.133/2021:

\begin{enumerate}
  \item \textbf{Jurídica}: ato constitutivo, estatuto ou contrato social em vigor;
  \item \textbf{Fiscal e Trabalhista}: certidões negativas de débitos federais, estaduais,
    municipais, FGTS e trabalhistas (TST);
  \item \textbf{Técnica}:
    \begin{itemize}
      \item Autorização de Funcionamento expedida pela ANVISA para distribuição de
        saneantes (para itens dos Lotes 01 a 05);
      \item Comprovação de capacidade logística fluvial para entrega em Borba/AM
        (carta ou contrato com empresa de navegação fluvial que opere no trecho
        Manaus--Borba);
      \item Atestado de capacidade técnica de fornecimento de objeto compatível em
        características e quantidades, conforme \art{67}, I;
    \end{itemize}
  \item \textbf{Econômico-financeira}: capital social ou patrimônio líquido mínimo
    de 10\% do valor do lote (\art{69}, I); certidão negativa de falência.
\end{enumerate}

\begin{warnbox}[Registro ANVISA --- Habilitação Técnica]
\small
Nos termos da Lei n.\textsuperscript{o}~6.360/1976 e da RDC ANVISA
n.\textsuperscript{o}~59/2021, \textbf{todos} os saneantes dos Lotes 01 a 05 devem
possuir registro ou notificação vigente na ANVISA. A entrega de produto com registro
vencido ou inexistente configura infração sanitária e enseja: (i) rescisão unilateral da
Ata; (ii) sanção ao fornecedor nos termos do \art{156}; e (iii) responsabilização do
agente público que atestou o recebimento.
\end{warnbox}

% ============================================================
\section{Declaração de Viabilidade da Contratação}
% ============================================================

\begin{legalbox}[Base Legal --- \art{18}, \S\,3.\textsuperscript{o},
  Lei n.\textsuperscript{o} 14.133/2021]
\small
O ETP deverá concluir com a \textbf{declaração de viabilidade ou inviabilidade} da
contratação. Concluindo pela viabilidade, deverá indicar a solução mais vantajosa.
\end{legalbox}

\begin{conclusabox}
\begin{center}
  \textbf{\color{navyblue}\large DECLARAÇÃO DE VIABILIDADE}
\end{center}

\medskip
Com base nas análises realizadas neste Estudo Técnico Preliminar, conclui-se pela
\textbf{VIABILIDADE} da contratação objeto deste documento, pelos seguintes fundamentos:

\begin{enumerate}
  \item A necessidade da contratação é \textbf{real, urgente e continuada},
    derivando de obrigação legal irrenunciável de manutenção da salubridade da rede
    de saúde que atende 34.869 habitantes;
  \item A solução via \textbf{Ata de Registro de Preços --- Pregão Eletrônico} é a
    mais adequada, flexível e economicamente vantajosa para a realidade logística
    fluvial de Borba/AM;
  \item A pesquisa de preços atesta a existência de \textbf{mercado competitivo}
    e preço de referência compatível com o orçamento disponível de
    \textbf{R\$~2.000.000,00};
  \item Os riscos foram identificados e suas mitigações planejadas,
    incluindo o \textbf{Estoque de Reserva Estratégica (ERE)} para o período
    de estiagem do Rio Madeira;
  \item A contratação atende aos princípios da \textbf{eficiência, economicidade e
    sustentabilidade} exigidos pela Lei n.\textsuperscript{o}~14.133/2021.
\end{enumerate}

\medskip
\textbf{Solução mais vantajosa}: Pregão Eletrônico, modo aberto, julgamento por menor
preço, dividido em 6 lotes, com instrumento de Ata de Registro de Preços, com vigência
de 12 meses prorrogáveis, e exigência de entrega CIF Borba/AM.
\end{conclusabox}

\vspace{1.5cm}

%% Assinaturas
\noindent
\begin{minipage}{0.47\linewidth}
\centering
\rule{6cm}{0.5pt}\\[4pt]
\textbf{[Nome e Cargo]}\\
Responsável pela Elaboração do ETP\\
Matrícula: \hspace{2cm}\\[2pt]
{\small Secretaria Municipal de Saúde de Borba/AM}
\end{minipage}
\hfill
\begin{minipage}{0.47\linewidth}
\centering
\rule{6cm}{0.5pt}\\[4pt]
\textbf{[Nome e Cargo]}\\
Autoridade Competente --- Aprovação\\
Matrícula: \hspace{2cm}\\[2pt]
{\small Secretaria Municipal de Saúde de Borba/AM}
\end{minipage}

\vspace{1.2cm}
\noindent\hfill{\small Borba/AM, \today}

\vspace{2cm}
\noindent\textcolor{navyblue!30}{\rule{\linewidth}{0.5pt}}

\noindent{\footnotesize\itshape
\textbf{Referências Normativas:}
Lei n.\textsuperscript{o}~14.133, de 1.\textsuperscript{o} de abril de 2021 (Nova Lei de
Licitações e Contratos Administrativos);
Decreto Federal n.\textsuperscript{o}~11.462/2023 (SRP);
IN SEGES/MGI n.\textsuperscript{o}~65/2021 (Pesquisa de Preços);
RDC ANVISA n.\textsuperscript{o}~15/2012 (Processamento de Artigos);
RDC ANVISA n.\textsuperscript{o}~59/2021 (Saneantes);
RDC ANVISA n.\textsuperscript{o}~216/2004 (Serviços de Alimentação);
NR-32 (MTE --- Segurança e Saúde no Trabalho em Serviços de Saúde);
LC n.\textsuperscript{o}~123/2006 (Estatuto da ME e EPP);
Lei n.\textsuperscript{o}~6.360/1976 (Vigilância Sanitária);
Lei n.\textsuperscript{o}~12.305/2010 (PNRS);
ABNT NBR 7.500/7.501 (Transporte e Armazenagem de Produtos Perigosos).}

\end{document}
